% Research draft: the relation and implementation remain to be specified and evaluated.
\section{Introduction}

Vehicular crowdsensing collects observations from vehicles that encounter conditions of interest while travelling through a road network. A task may request a report only when its observation falls within a designated region and time interval. An operator accepting such reports must establish that each submission satisfies the task policy while protecting information that could reveal the participant's identity or movement. Prior work on privacy-preserving crowdsensing has studied anonymous participation, encrypted reports and duplicate-submission detection in related settings~\cite{Ni2020SPOON,Xu2021TaskMatching}.

The focus here is the acceptance of one report for one task. We consider a verifier that receives a report and associated evidence, but need not learn the participant's identity or the exact location and time used to satisfy the policy. The desired decision is whether an authorised credential holder submitted a report associated with an eligible spatiotemporal event and whether the same credential has already been used for that task. This per-report decision complements coverage auditing, which concerns the accumulated coverage of regions over time~\cite{He2026TSTAudit}.

These requirements have to refer to the same submission. A participant might possess a valid credential, a valid location-related attestation and a separately prepared report, yet their independent validity does not itself establish that the attestation and the report refer to the same event or intended credential holder. We therefore study a relation that joins a task identifier and policy, an authorised credential, an issuer's event evidence, and a commitment to the submitted report. The evidence is intended to be bound to the same credential holder and event as the report commitment. The exact signed fields, issuer trust model and proof statement must be specified before the binding guarantee can be assessed.

The planned construction uses a task-scoped deterministic tag derived from credential-held secret material to support duplicate detection without directly disclosing the credential holder. Equal tags within a task reveal repeated use of the same credential; different task scopes are intended to yield unlinkable tags, subject to the cryptographic construction and other visible metadata. This provides a per-credential limit, not necessarily a per-person or per-vehicle limit: the latter depends on registration and credential-issuance policy. Region and time membership are intended to be verified through a privacy-preserving proof over the event evidence and the task policy. The visibility of task identifiers, policy parameters, report data and network metadata must be stated in the protocol model.

The proposed study has three objectives. First, it will define the acceptance relation and the roles of the credential issuer, event-evidence issuer, participant and verifier. Second, it will specify how report, event, issuer and task references are bound in that relation, alongside a scoped duplicate-detection rule. Third, it will evaluate the resulting construction using positive cases and adversarial substitution or replay cases, with performance measured for the implemented operations. These are research objectives; the relation, security claims and results remain to be established by the completed protocol and evaluation.

The event issuer and any trusted acquisition component are explicit assumptions. A signed or committed report can establish integrity and association with an attestation under those assumptions; it does not by itself establish the physical truth of the observation. In particular, a fresh challenge can make the resulting evidence fresh without proving that the underlying sensor reading was acquired after the challenge~\cite{Birkholz2023RATS}. A trusted capture path or another justified acquisition mechanism is needed for any stronger capture-time claim. Likewise, signature and proof verification alone do not establish resistance to a real-time relay.

\section{Related Work}

\paragraph{Privacy-preserving crowdsensing and duplicate detection.}
Ni et al.\ developed SPOON for privacy-preserving task allocation and reporting, combining anonymous credentials, encrypted sensing reports and mechanisms to detect duplicate reporting~\cite{Ni2020SPOON}. Xu et al.\ study privacy-preserving task matching and multiple-submission detection, including a restriction of one submission per qualified worker per task~\cite{Xu2021TaskMatching}. Zhu et al.\ consider privacy-preserving hybrid sensing and detection of duplicate reports in a smart-parking application~\cite{Zhu2022HybridSensing}. These studies establish that anonymous participation and scoped duplicate control are established research directions. The question considered here is the explicit joint relation between the task policy, credential, issuer evidence and report commitment; whether a particular construction advances beyond prior work depends on the complete protocol and a detailed comparison.

\paragraph{Spatiotemporal coverage auditing.}
He et al.\ propose TST-Audit for trustworthy spatiotemporal coverage auditing in vehicular crowdsensing~\cite{He2026TSTAudit}. Its design includes a microproof interface for mapping a hidden location and time to a public grid cell and epoch, credential-based uniqueness within that scope, and authenticated aggregate coverage. The published coverage-record interface is centred on a grid, epoch, credential-derived value and microproof. The present research instead takes an individual sensing report and its task-specific acceptance policy as the decision object, and investigates an explicit binding to report content and issuer evidence. TST-Audit's microproof is an abstract interface, so this distinction does not imply that every possible instantiation of that interface lacks report binding.

\paragraph{Location proofs and evidence provenance.}
Proof-of-location systems have considered multiple witnesses, distance constraints and time consistency~\cite{Brito2025DecentralizedPoL}. Brito et al.\ also investigate a witnessing-zone architecture for provenance of captured content, where multiple observers validate physical events~\cite{Brito2026ContentProvenance}. Their work is directly relevant to the link between content and its physical context; event provenance should therefore not be presented as a new problem. Bogdanov et al.\ study zero-knowledge proofs of vehicle-policy compliance using location information supplied by a trusted device~\cite{Bogdanov2025VehiclePolicy}, while Ernstberger et al.\ address zero-knowledge proofs of location properties using accurate floating-point arithmetic~\cite{Ernstberger2025ZKLP}. These works motivate careful treatment of both the policy predicate and the source of location evidence. Their precise guarantees must be compared against the final statement and trust assumptions of the proposed protocol.

\paragraph{Sensor data authenticity.}
IoT Notary studies attestable capture and logging of sensor data~\cite{Panwar2022IoTNotary}, while BlockSense addresses trustworthy mobile crowdsensing and data quality~\cite{Huang2024BlockSense}. These approaches show why authenticity and quality cannot be inferred solely from a successful policy proof. The planned relation concerns the acceptance of a report under a stated issuer and acquisition model; any claim about when or how the sensor reading was physically obtained requires evidence from that model.
